% TOP_PROBLEM: {"id": 30007038, "title": "Kähler fourfold abundance after nonvanishing", "category_id": 5, "category": {"id": 5, "name": "algebraic_geometry", "display_name": "Algebraic Geometry", "description": "Geometric objects defined by polynomial equations.", "slug": "algebraic-geometry", "order_index": 5, "created_at": "2026-07-31T15:26:25.671Z"}, "status": "open"}
% FIRST_POSTED: 2026-10-06
% ATTEMPT_STATUS: solved
% OPENAI_CLAIM: {"schema_version":1,"source_repo":"https://github.com/openai/math","source_commit":"adc7f1241b42e322a6451854ab7e4b4c146bf78a","release_date":"2026-10-06","match":"full","resolution":"proved","independently_reviewed":false,"second_pass":"agreed","families":[{"family":"034","title":"Log abundance for compact Kähler spaces under logarithmic Iitaka subadditivity","subject":"Algebraic and complex geometry","manuscripts":[{"title":"Abundance after nonvanishing for compact Kähler fourfolds","pdf_path":"preprints/Abundance-after-nonvanishing-for-compact-Kahler-fourfolds-September-27-2026/paper.pdf","readme_path":"preprints/Abundance-after-nonvanishing-for-compact-Kahler-fourfolds-September-27-2026/README.md","date":"2026-09-27","theorem_ref":"thm:main (sections/introduction.tex:23)","statement_sources":[{"path":"preprints/Abundance-after-nonvanishing-for-compact-Kahler-fourfolds-September-27-2026/build/sections/introduction.tex","line":23,"label":"thm:main"}]}],"lean":null,"reasoning_trace":null}]}
\documentclass[11pt]{article}
\usepackage[margin=25mm]{geometry}
\usepackage{fontspec}
\setmainfont{Latin Modern Roman}
\usepackage{amsmath,amssymb,mathtools}
\usepackage{unicode-math}
\setmathfont{Latin Modern Math}
\usepackage{enumitem}
\usepackage{xurl,hyperref,bookmark}
\hypersetup{colorlinks=true,urlcolor=blue}
\setcounter{secnumdepth}{0}
\setlength{\parindent}{0pt}
\setlength{\parskip}{5pt}
\setlength{\emergencystretch}{3em}
\begin{document}
\section*{Kähler fourfold abundance after nonvanishing}

\subsection{OpenAI's claim}
A nef rational adjoint on a compact Kähler klt fourfold is claimed to be semiample once a nonzero section exists, including triviality of a multiple in Iitaka dimension zero. The cited manuscript argues for that result; its relation to this problem is specified below.

\subsection{Mathematical statement of the claimed result}
\paragraph{Mathematical claim 1.} Let $X$ be a compact complex space, required to be Kähler, normal and connected, of dimension four and let $\Delta\ge0$ be rational, with $(X,\Delta)$ Kawamata log terminal and actual adjoint $D=K_X+\Delta$ rational Cartier. Assume $D$ is analytically nef: after fixing a Cartier multiple and a Kähler form, its curvature has smooth representatives bounded below by $-\varepsilon$ times that form for every $\varepsilon>0$. Assume also $\kappa(X,D)\ge0$, meaning a nontrivial section exists for at least one positive Cartier multiple. OpenAI claims
\[
\exists m>0:\ mD\text{ Cartier},\qquad
H^0(X,\mathcal O_X(mD))\otimes\mathcal O_X\twoheadrightarrow\mathcal O_X(mD).
\]
If $\kappa(X,D)=0$, the multiple may be chosen with $\mathcal O_X(mD)\simeq\mathcal O_X$. This theorem requires nonvanishing as a hypothesis.

\subsection{Outline of the argument}
Manuscript 1 (family 034, 2026-09-27): Contractions of threefold null loci and analytic thickenings prepare the boundary components for a fourfold birational reduction. The comparison on floors and lower strata preserves the actual holomorphic line identity, allowing boundary-generated sections to descend to the nef adjoint.

\subsection{Scope relative to this problem}
Manuscript 1 (family 034, 2026-09-27): The problem selects the mathematical statement summarized here, including its stated hypotheses and parameter range. The cited theorem gives a claimed proof of that assertion in full.

\subsection{Formal verification}
Family 034 has no Lean scope page in this snapshot; no formalization of these manuscript results is listed.

This site has not executed the Lean files or checked their dependencies. These descriptions report the scope of the linked declarations, rather than a new verification of the proofs.

\subsection{Status caveat}
OpenAI attributes these results to an internal model. Its release notice cautions that results lacking full formalization can contain mistakes. This site has not independently assessed the proofs; a solved label records the scope of the reported claim.

\subsection{Sources}
\begin{itemize}
\item OpenAI, \emph{Abundance after nonvanishing for compact Kähler fourfolds}, 2026-09-27; citation key \texttt{OAI:Abundance-after-nonvanishing-for-compact-Kahler-fourfolds-September-27-2026}. \href{https://github.com/openai/math/blob/adc7f1241b42e322a6451854ab7e4b4c146bf78a/preprints/Abundance-after-nonvanishing-for-compact-Kahler-fourfolds-September-27-2026/paper.pdf}{PDF (GitHub)}; \href{https://github.com/openai/math/raw/adc7f1241b42e322a6451854ab7e4b4c146bf78a/preprints/Abundance-after-nonvanishing-for-compact-Kahler-fourfolds-September-27-2026/paper.pdf}{PDF (direct)}; \href{https://github.com/openai/math/blob/adc7f1241b42e322a6451854ab7e4b4c146bf78a/preprints/Abundance-after-nonvanishing-for-compact-Kahler-fourfolds-September-27-2026/README.md}{Manuscript page and citation}; \href{https://github.com/openai/math/blob/main/preprints/Abundance-after-nonvanishing-for-compact-Kahler-fourfolds-September-27-2026/paper.pdf}{Latest version}; Statement source: \href{https://github.com/openai/math/blob/adc7f1241b42e322a6451854ab7e4b4c146bf78a/preprints/Abundance-after-nonvanishing-for-compact-Kahler-fourfolds-September-27-2026/build/sections/introduction.tex\#L23}{thm:main}.
\item \href{https://github.com/openai/math/blob/adc7f1241b42e322a6451854ab7e4b4c146bf78a/CONTENTS.md}{OpenAI catalogue entry}.
\end{itemize}
\end{document}
